\section{Conclusion}
\label{sec:conclusion}

We compared a plain ViT (\medsam{}), a hybrid (\litemedsam{}) and three CNNs on the thin vessels of three retinal datasets, DRIVE, HRF and FIVES, with VessMap as a check in a second imaging domain.

Fine-tuned from its pre-trained weights and evaluated with the upscaled protocol, \medsam{} had an $\thinrecall$ not significantly different from that of the best CNN on all four datasets, and a significantly lower $P_{\mathrm{thin}}$ on DRIVE and VessMap. For \medsam{} in these conditions, the results on these datasets do not support the expectation that a plain ViT misses thin vessels.

\medsam{} reached this as a model of the most expensive class by parameter count. A compact plain ViT trained from scratch, in the class of \litemedsam{} and \convunext{}, did not keep this performance: against the best CNN, its $\thinrecall$ was significantly lower on DRIVE and FIVES, and its $P_{\mathrm{thin}}$ and Dice on all four datasets. The hybrid \litemedsam{}, in the same class, had the highest mean $\thinrecall$ on every dataset, but a lower $P_{\mathrm{thin}}$ than the best CNN, significantly on three datasets. Where cost matters, the smallest CNN remains a candidate on small images: \lwnet{}, with 0.07~M parameters, had a Dice and an $\thinrecall$ not significantly different from those of fine-tuned \medsam{} on DRIVE and VessMap, and on DRIVE a significantly higher $P_{\mathrm{thin}}$ (Table~\ref{tab:pvalues}), at a small fraction of its cost per image (Fig.~\ref{fig:cost}).

Pre-training raised the Dice of \medsam{} significantly on every dataset; without it, \medsam{} fell significantly below the best CNN in $\thinrecall$ on three datasets. The hybrid \litemedsam{} depended less on pre-training for Dice than for $\thinrecall$. The Dice gain of \medsam{} was smallest on FIVES, the dataset with the most training images, while the gains of \litemedsam{} did not follow this order; four datasets that also differ in resolution, domain and upscaling cannot isolate the number of images. On HRF and FIVES, fine-tuned \medsam{} had a significantly higher Dice than the best CNN, with no significant loss in $\thinrecall$ or $P_{\mathrm{thin}}$.

On DRIVE and VessMap, whose images are smaller than the window of \medsam{}, padding the images instead of upscaling them to the window lowered its $\thinrecall$ below the range of the CNNs at the same cost. For \medsam{} on these two datasets, the thin-vessel results depend on evaluating it at the scale it was trained on, and a study of a plain ViT on images smaller than its window should state that scale.

Each configuration was trained once, so the tests generalize over test images and not over training runs, and the comparison across architectures mixes initializations and input scales. Future work should repeat the comparison over several training runs and on more datasets, train the CNNs at the input scale and window of \medsam{}, train the plain ViT from scratch with a recipe designed for it, test the patch size of a plain ViT with the other hyperparameters fixed, and measure inference latency.

\section*{Code availability}
The code for training, the hyperparameter search and the figures is available at \url{https://github.com/CarlosFontaneli/thin-vessel-sam-cnn}.
